\section{Asymptotics of the generalized hypergeometric function}\label{Sect. 4}
Our derivation in Section~\ref{Sect. 3} depends on a rough estimate of $_2F_2$ for large $-n$. In this section, Nagel's approach is also used to explicitly establish the asymptotic behavior of $_2F_2$ for large parameters. We also present the asymptotic behavior of $_pF_q$ for large $-n$.

\subsection[Asymptotics of \_2F\_2]{Asymptotics of $\boldsymbol{_2F_2}$}
Let us examine the complete asymptotic expansions of $_2F_2$ for large parameters.

\begin{Theorem}\label{Thm: 2F2 large parameter--lambda}
	Let $\delta\in (0,\pi)$, $c\notin\ZZ_{\leqslant 0}$ and $z\ne 0$. Then for any positive integer $N$,
	\begin{equation}\label{2F2 large parameter--lambda}
		{}_2F_2\biggl[\begin{matrix}
			a,b+\lambda\\
			c,d+\lambda
		\end{matrix};z\biggr]=\sum_{k=0}^{N-1}\frac{(a)_k (d-b )_k}{(c)_k (d+\lambda )_k}\frac{ (-z )^k}{k!}{}_1F_1\biggl[\begin{matrix}
			a+k\\
			c+k
		\end{matrix};z\biggr]+\mo \bigl(\lambda^{-N} \bigr),
	\end{equation}
	where $\lambda\to\infty$ in the sector $|{\arg}(\lambda+d)|\leqslant \pi-\delta$.
\end{Theorem}

\begin{proof}
	If denoting
	\[v(z):={}_1F_1\biggl[\begin{matrix}
		a\\
		c
	\end{matrix};z\biggr],\qquad d_k:=\frac{(a)_k}{(c)_k k!},\]
	we can obtain
	\[{}_2F_2\biggl[\begin{matrix}
		a,b+\lambda\\
		c,d+\lambda
	\end{matrix};z\biggr]=\sum_{k=0}^{\infty}\frac{ (\lambda+b )_k}{ (\lambda+d )_k}d_k z^k.\]
	Note that the $k$-th derivative of $v(z)$ is given by \cite[p.~405, equation~(16.3.1)]{NIST-Handbook}
	\[v^{(k)}(z)=\frac{(a)_k}{(c)_k}
	{}_1F_1\biggl[\begin{matrix}
		a+k\\
		c+k
	\end{matrix};z\biggr].\]
	Applying Fields' result \cite[Theorem 3]{Fields-1967} yields
	\begin{gather}\label{2F2 transformation}
		{}_2F_2\biggl[\begin{matrix}
			a,b+\lambda\\
			c,d+\lambda
		\end{matrix};z\biggr]=\sum_{k=0}^{\infty}\frac{(d-b)_k}{(d+\lambda)_k}\frac{(-z)^k}{k!}v^{(k)}(z)=\sum_{k=0}^{\infty}\frac{(a)_k(d-b)_k}{(c)_k(d+\lambda)_k}\frac{(-z)^k}{k!}{}_1F_1\biggl[\begin{matrix}
			a+k\\
			c+k
		\end{matrix};z\biggr]
	\end{gather}
	with $\lambda+d\notin\ZZ_{\leqslant 0}$, and also gives the asymptotic expansion \eqref{2F2 large parameter--lambda}.
\end{proof}

\begin{Remark}
	~
	\begin{itemize}\itemsep=0pt
		\item[(1)] The series (\ref{2F2 transformation}) can be reformulated as follows:
		\begin{equation}\label{2F2 transformation reformulated}
			{}_2F_2\biggl[\begin{matrix}
				a,b\\
				c,d
			\end{matrix};z\biggr]=\sum_{k=0}^{\infty}\frac{(a)_k(d-b)_k}{(c)_k(d)_k}\frac{(-z)^k}{k!}{}_1F_1\biggl[\begin{matrix}
				a+k\\
				c+k
			\end{matrix};z\biggr],
		\end{equation}
		which is a specialization of Luke's result \cite[Section~9.1, equation~(34)]{Luke-1969}
		\begin{align}
					&{}_{p+1}F_{q+1}\biggl[\begin{matrix}
						a_1,\dots,a_p,b\\
						c_1,\dots,c_q,d
					\end{matrix};z\biggr]{}\nonumber\\
 & \qquad{}= \sum_{k=0}^{\infty}\frac{ (a_1 )_k\cdots (a_p )_k}{ (c_1 )_k\cdots (c_q )_k}\frac{(d-b)_k}{(d)_k}
 \frac{(-z)^k}{k!}{}_{p}F_{q}\biggl[\begin{matrix}
						a_1+k,\dots,a_p+k\\
						c_1+k,\dots,c_q+k
					\end{matrix};z\biggr],\label{pFp transformation}
\end{align}
			where $p\leqslant q$ with $z\in\CC$, or $p=q+1$ with $\operatorname{Re}(z)<\frac{1}{2}$.
		
		\item[(2)] Jur\v{s}\.enas \cite[Section~4.1, p.~69]{Jursenas-2014} mentioned an asymptotic formula
		\begin{equation}\label{Jursenas expansion-1}
			{}_2F_2\biggl[\begin{matrix}
				\alpha+1,\beta+\nu\\
				\delta,\beta+1+\nu
			\end{matrix};-\frac{1}{z}\biggr]\sim {}_1F_1\biggl[\begin{matrix}
				\alpha+1\\
				\delta
			\end{matrix};-\frac{1}{z}\biggr],\qquad \nu\in\ZZ_{\geqslant 0},\quad\nu\to \infty.
		\end{equation}
		without proof, whereas our Theorem \ref{Thm: 2F2 large parameter--lambda} gives a full asymptotic expansion of (\ref{Jursenas expansion-1}).
	\end{itemize}
\end{Remark}

The asymptotics of $_2F_2$ for $\lambda=-n$ is given below.
\begin{Theorem}\label{Thm: 2F2 large parameter--n}
	Let $c\notin\ZZ_{\leqslant 0}$, $d\notin\ZZ$, $b-d\notin\ZZ_{\geqslant 0}$ and $z\ne 0$. Then for any positive integer $N$,
	\begin{equation}\label{2F2 large parameter--n}
		{}_2F_2\biggl[\begin{matrix}
			a,b-n\\
			c,d-n
		\end{matrix};z\biggr]=\sum_{k=0}^{N-1}\frac{(a)_k(d-b)_k}{(c)_k(d-n)_k}\frac{(-z)^k}{k!}{}_1F_1\biggl[\begin{matrix}
			a+k\\
			c+k
		\end{matrix};z\biggr]+\mo\big(n^{-N}\big),
	\end{equation}
	where $n\to +\infty$ through integer values.
\end{Theorem}

\begin{proof}
	We shall follow Nagel's approach. For nonnegative integers $n$ and $\ell$, write
	\[F(z):={}_2F_2\biggl[\begin{matrix}
		a,b-n\\
		c,d-n
	\end{matrix};z\biggr],\qquad a_{\ell}(n):=\frac{(-1)^{\ell}}{(d-n)_{\ell}},\qquad g_{\ell}:=\frac{(a)_{\ell}(d-b)_{\ell}}{(c)_{\ell}\ell!}.\]
	Clearly, $\left|g_\ell\right|\leqslant C\ell^{\operatorname{Re}(\alpha)}$, where $\alpha:=a+d-b-c-1$. Moreover, \eqref{2F2 transformation reformulated} suggests that
	\begin{equation}\label{R(z) definition}
		R(z):=F(z)-G_0(z)=\sum_{\ell=1}^{\infty}a_\ell(n)g_\ell G_\ell(z)z^{\ell},
	\end{equation}
	where $G_\ell(z)$ is given by \eqref{G_l definition}.
	
	Take $n$ large, let $m=\lceil \log n\rceil$ and divide the series \eqref{R(z) definition} into five parts:
	\begin{equation}\label{five sums}
		R(z)=\sum_{1}^{N}+\sum_{N+1}^m+\sum_{m+1}^{n/2}+\sum_{n/2+1}^n+\sum_{n+1}^{\infty}=:S_0+S_1+S_2+S_3+S_4,
	\end{equation}
	where $N\in\ZZ_{>0}$ is chosen so that $\operatorname{Re}(c+N+1)>0$. Therefore, the inequality \eqref{G_l inequality} holds.
	
	Let us derive estimates for the sums in \eqref{five sums}.
	
	\textbf{Case 1.} $1\leqslant \ell\leqslant N$. Now $a_\ell(n)\sim n^{-\ell}$. Thus
	\[S_0=\sum_{\ell=1}^{N-1}a_\ell(n)g_\ell G_\ell(z)z^{\ell}+\mo\big(n^{-N}\big).\]
	
	\textbf{Case 2.} $N+1\leqslant\ell\leqslant m$. Now both $n$ and $n-\ell$ are large. The use of the identity $(z)_\ell=(-1)^{\ell}(1-z-\ell)_{\ell}$ gives
	\begin{equation}\label{a_l(n) identity}
		a_\ell(n)=\frac{1}{(1-d+n-\ell)_\ell}=\frac{\Gamma(1-d+n-\ell)}{\Gamma(1-d+n)}.
	\end{equation}
	By Stirling's formula, we get $a_\ell(n)\sim n^{-\ell}$ and thus
	\[|S_1|\leqslant C\mathrm{e}^{\gamma|z|}\sum_{\ell=N+1}^{m}\ell^{\operatorname{Re}(\alpha)}\biggl(\frac{|z|}{n}\biggr)^{\ell}.\]
	For $n$ large, $\ell^{\operatorname{Re}(\alpha)}\big(\frac{|z|}{n}\big)^{\ell}$ decreases with respect to $\ell\geqslant N+1$ and then
	\[|S_1|\leqslant C\mathrm{e}^{\gamma|z|}m(N+1)^{\operatorname{Re}(\alpha)}\biggl(\frac{|z|}{n}\biggr)^{N+1}=\mo\biggl(\frac{\log n}{n^{N+1}}\biggr).\]
	
	\textbf{Case 3.} $m+1\leqslant \ell\leqslant \frac{n}{2}$. Stirling's formula shows that
	\[a_\ell(n)=\frac{\Gamma(1-d+n-\ell)}{\Gamma(1-d+n)}\sim \biggl(1-\frac{\ell}{n}\biggr)^{-d}\frac{(n-\ell)!}{n!}.\]
	It follows from $\frac{1}{2}\leqslant 1-\frac{\ell}{n}<1$ that
	\[|a_\ell(n)|\leqslant C \frac{(n-\ell)!}{n!}=\frac{C}{(n-\ell+1)\cdots(n-1)n}\leqslant C\bigg(\frac{2}{n}\bigg)^{\ell}.\]
	As in Case 2, the monotonicity gives
	\[|S_2|\leqslant C\mathrm{e}^{\gamma|z|}\sum_{\ell=m+1}^{n/2}\ell^{\operatorname{Re}(\alpha)}\bigg(\frac{2|z|}{n}\bigg)^{\ell}\leqslant C\mathrm{e}^{\gamma|z|}nm^{\operatorname{Re}(\alpha)}\bigg(\frac{2|z|}{n}\bigg)^{m+1}=\mo\big(n^{-\frac{1}{2}\log n}\big).\]
	
	\textbf{Case 4.} $\frac{n}{2}+1\leqslant \ell\leqslant n$. Choose the least number $r\in\ZZ_{\geqslant 0}$ so that
	\[|(1-d+n-\ell)+r|\geqslant 4|z|>0.\]
	Note that $\frac{1}{|z+n|}$ is bounded uniformly for $n\geqslant 1$. Using \eqref{a_l(n) identity} gives
	\[|a_{\ell}(n)|=\prod_{j=0}^{\ell-1}\big|(1-d+n-\ell)+j\big|^{-1}\leqslant C(4|z|)^{r-\ell}.\]
	Thus
	\[|S_3|\leqslant C\mathrm{e}^{\gamma|z|}(4|z|)^r\sum_{\ell=n/2+1}^{n}\ell^{\operatorname{Re}(\alpha)}4^{-\ell}=\mo\big(n^{\operatorname{Re}(\alpha)+1}2^{-n}\big).\]
	
	\textbf{Case 5.} $\ell\geqslant n+1$. Now we can obtain $|a_\ell(n)|\leqslant C\frac{n^{\operatorname{Re}(d)}(\ell-n)^{1-\operatorname{Re}(d)}}{n!(\ell-n)!}$ from
	\[a_\ell(n)=\frac{(-1)^{\ell-n}}{(1-d)_n(d)_{\ell-n}}=\frac{(1)_n(1)_{\ell-n}}{(1-d)_n(d)_{\ell-n}}\frac{(-1)^{\ell-n}}{n!(\ell-n)!}.\]
	It follows that
	\begin{align*}
\begin{aligned}
		|S_4|&{} \leqslant C\mathrm{e}^{\gamma|z|}\frac{n^{\operatorname{Re}(d)}}{n!}\sum_{\ell=n+1}^{\infty}\ell^{\operatorname{Re}(\alpha)}(\ell-n)^{1-\operatorname{Re}(d)}\frac{|z|^{\ell}}{(\ell-n)!}\\
		&{} =C\mathrm{e}^{\gamma|z|}n^{\operatorname{Re}(d)}\frac{|z|^n}{n!}\sum_{j=1}^{\infty}(j+n)^{\operatorname{Re}(\alpha)}j^{1-\operatorname{Re}(d)}\frac{|z|^j}{j!}.
\end{aligned}
	\end{align*}
	Note that $\max\{n,j\}\leqslant j+n\leqslant 2\cdot\max\{n,j\}$. Then for $p,q\in\RR$, we have
	\begin{align*}
		\sum_{j=1}^{\infty}(j+n)^p j^q\frac{|z|^j}{j!}&{}\leqslant Cn^p\sum_{j=1}^{n}j^q\frac{|z|^j}{j!}+C\sum_{j=n+1}^{\infty}j^{p+q}\frac{|z|^j}{j!}\\
		&{} \leqslant Cn^{\max\{0,p\}}\sum_{j=1}^{\infty}j^{\max\{q,p+q\}}\frac{|z|^j}{j!},
	\end{align*}
	which shows that
	\[|S_4|\leqslant C\mathrm{e}^{\gamma|z|}n^{\max\{\operatorname{Re}(d),\operatorname{Re}(\alpha+d)\}}\frac{|z|^n}{n!}\sum_{j=1}^{\infty}j^{1-\operatorname{Re}(d)+\max\{0,\operatorname{Re}(\alpha)\}}\frac{|z|^j}{j!}.\]
	
	Now the asymptotic expansion \eqref{2F2 large parameter--n} follows from the estimates above.
\end{proof}
